% P8 related work
\section{Related Work}
\label{sec:p8-related}

\paragraph{Gradient boosting on tabular data.} XGBoost
\citep{chen2016xgboost} remains the workhorse of applied tabular
classification. \citet{grinsztajn2022tree} show systematically that tree-based
models still outperform deep learning on medium-sized tabular benchmarks,
attributing the gap to trees' robustness to uninformative features and to the
non-smooth, non-rotation-invariant structure of tabular targets;
\citet{shwartzziv2021tabular} reach the same conclusion across deep tabular
architectures. Our scorer result is a single applied data point consistent
with this literature, with the added operational arguments (latency, cost,
calibration \citep{guo2017calibration}, injection surface
\citep{greshake2023injection}) that the benchmarking papers do not address.

\paragraph{LLMs for tabular data.} TabLLM \citep{hegselmann2023tabllm}
established the serialize-and-prompt recipe and showed it is most competitive
in very-few-shot regimes, degrading relative to boosted trees as labeled data
grows---exactly the pattern our head-to-head reproduces at $40{,}000$ labels,
and exactly why our taxonomy assigns the LLM the cold-start seat
(\secref{sec:p8-gap-coldstart}) rather than the scorer seat.

\paragraph{Fraud detection on transaction data.} \citet{dalpozzolo2018realistic}
give the canonical treatment of realistic credit-card fraud modeling---class
imbalance, verification latency, and concept drift---and motivate the
anonymized-feature, heavily imbalanced regime our synthetic benchmark
imitates. Our contribution is orthogonal to this line: we do not propose a
better detector but a division of labor around an already-strong one.

\paragraph{LLM agents for financial crime.} \citet{pirmorad2025icl_aml} shows
that LLMs prompted with serialized financial-graph neighborhoods can perform
analyst-style money-laundering triage few-shot, with coherent red-flag
justifications; \citet{naik2025coinvestigator} builds an agentic framework
specifically for drafting AML compliance narratives (SARs) with
human-in-the-loop review. These systems instantiate, respectively, our
cold-start and scribe/triage seats; our contribution is to place them in an
explicit architecture where the real-time scorer remains a tree.

\paragraph{VLMs for document and image fraud.} FakeShield
\citep{xu2024fakeshield} demonstrates explainable image-forgery detection and
localization with multimodal LLMs; Forensics-Bench \citep{wang2025forensicsbench}
benchmarks large vision--language models across 112 forgery types and finds
substantial headroom. Together they support both halves of our sensor claim:
VLMs uniquely extend the input space to raw document evidence, and their
current reliability justifies confining them to feature extraction with
downstream tree scoring and human review rather than autonomous verdicts.

\paragraph{Program context.} The parent reproducibility program measures
RL post-training claims capability-by-capability rather than by their labels;
the present side-probe applies the identical discipline to the marketing label
``AI'' in fraud detection. The methodological through-line is decomposition:
a bundled claim is replaced by per-capability measurements, and the deployment
decision is made seat by seat.
